% =====================================================================
% Fiche de revision - Python Scripting - Lesson 1
% Compilation : pdfLaTeX (Overleaf), aucune option particuliere requise.
% =====================================================================

\documentclass[a4paper,10pt]{article}

\usepackage[utf8]{inputenc}
\usepackage[T1]{fontenc}
\usepackage[french]{babel}
\usepackage[a4paper,margin=1.3cm]{geometry}
\usepackage{xcolor}
\usepackage{multicol}
\usepackage{booktabs}
\usepackage{array}
\usepackage{tabularx}
\usepackage{enumitem}
\usepackage{listings}
\usepackage{titlesec}
\usepackage{tcolorbox}

% ---------------------------------------------------------------- couleurs
\definecolor{accent}{HTML}{1F5FA8}
\definecolor{codefg}{HTML}{A03030}
\definecolor{codebg}{HTML}{F4F6F8}
\definecolor{kw}{HTML}{0B5FA5}
\definecolor{cmt}{HTML}{6C8A54}
\definecolor{strcol}{HTML}{2E7D32}
\definecolor{danger}{HTML}{B3261E}

% ---------------------------------------------------------------- titres
\titleformat{\section}
  {\normalfont\large\bfseries\color{accent}}{}{0pt}{}
  [\vspace{-0.75em}{\color{accent}\rule{\linewidth}{0.7pt}}]
\titlespacing*{\section}{0pt}{1.4ex}{0.7ex}

\titleformat{\subsection}[runin]
  {\normalfont\bfseries}{}{0pt}{}[\;---\;]
\titlespacing*{\subsection}{0pt}{0.8ex}{0.4em}

% ---------------------------------------------------------------- code
\lstset{
  language=Python,
  basicstyle=\ttfamily\footnotesize,
  keywordstyle=\color{kw}\bfseries,
  commentstyle=\color{cmt}\itshape,
  stringstyle=\color{strcol},
  showstringspaces=false,
  columns=fullflexible,
  keepspaces=true,
  backgroundcolor=\color{codebg},
  frame=none,
  xleftmargin=0.3em,
  aboveskip=0.55em,
  belowskip=0.55em,
  breaklines=true,
}

% code en ligne
\newcommand{\co}[1]{\texttt{\textcolor{codefg}{#1}}}

% ---------------------------------------------------------------- boites
\newtcolorbox{piegebox}[1]{
  colback=red!3, colframe=danger, boxrule=0.7pt,
  left=5pt, right=5pt, top=4pt, bottom=4pt,
  title=#1, fonttitle=\bfseries\small, coltitle=white
}

\newtcolorbox{astucebox}[1]{
  colback=accent!5, colframe=accent, boxrule=0.7pt,
  left=5pt, right=5pt, top=4pt, bottom=4pt,
  title=#1, fonttitle=\bfseries\small, coltitle=white
}

\setlist[itemize]{leftmargin=1.1em, itemsep=0.15em, topsep=0.25em, parsep=0pt}
\setlength{\columnsep}{1.1em}
\setlength{\parindent}{0pt}
\renewcommand{\arraystretch}{1.15}

\pagestyle{empty}

% =====================================================================
\begin{document}

% Desactive les espaces automatiques du francais avant : ; ! ?
% pour que le code en ligne (ex. "if a > 0:") ne soit pas deforme.
\shorthandoff{;:!?}

\small

\begin{center}
  {\LARGE\bfseries\color{accent} Python Scripting --- Lesson 1}\\[0.3em]
  {\normalsize Fiche de révision QCM \quad$\cdot$\quad les points qui tombent}
\end{center}
\vspace{0.4em}

% =====================================================================
\begin{piegebox}{Les 16 réponses pièges à connaître par c\oe ur}
\begin{minipage}[t]{0.485\linewidth}
\begin{tabularx}{\linewidth}{@{}l>{\raggedright\arraybackslash}X@{}}
\toprule
\textbf{Expression} & \textbf{Résultat} \\
\midrule
\co{int(3.7)}      & \co{3} --- tronque, n'arrondit pas \\
\co{int(-3.7)}     & \co{-3} --- tronque vers zéro \\
\co{10 / 2}        & \co{5.0} --- \co{/} rend un \co{float} \\
\co{10 // 3}       & \co{3} --- division entière \\
\co{-7 // 2}       & \co{-4} --- arrondi vers le bas \\
\co{10 \% 3}       & \co{1} --- le reste \\
\co{2 ** 3 ** 2}   & \co{512} --- de droite à gauche \\
\co{bool("False")} & \co{True} --- chaîne non vide \\
\bottomrule
\end{tabularx}
\end{minipage}\hfill
\begin{minipage}[t]{0.485\linewidth}
\begin{tabularx}{\linewidth}{@{}l>{\raggedright\arraybackslash}X@{}}
\toprule
\textbf{Expression} & \textbf{Résultat} \\
\midrule
\co{lst.sort()}        & \co{None} --- trie sur place \\
\co{mot.replace(...)}  & nouvelle chaîne, \co{mot} intact \\
\co{lst.append([4,5])} & ajoute \textbf{une} sous-liste \\
\co{[[0]*3]*3}         & 3 fois \textbf{la même} ligne \\
\co{zip} tailles $\neq$ & s'arrête à la plus courte \\
\co{for...else} + \co{break} & le \co{else} est ignoré \\
\co{d.get("z")}        & \co{None}, pas de \co{KeyError} \\
\co{(1)} vs \co{(1,)}  & \co{int} vs \co{tuple} \\
\bottomrule
\end{tabularx}
\end{minipage}
\end{piegebox}

\vspace{0.5em}
\begin{multicols}{2}

% =====================================================================
\section{Types et conversions}

Le typage est \textbf{dynamique} : une variable peut changer de type. Les
annotations (\co{age: int = 25}) sont de simples indications, \textbf{jamais
vérifiées} à l'exécution.

\begin{lstlisting}
int(3.7)        # 3    (tronque)
int("42")       # 42
int("1010", 2)  # 10   (lu en binaire)
int(True)       # 1
float(5)        # 5.0
str([])         # "[]"
\end{lstlisting}

\subsection{Les deux erreurs de conversion}
\co{ValueError} quand la conversion est impossible pour \emph{cette valeur}
(\co{int("abc")}) ; \co{TypeError} quand le \emph{type} est incompatible
(\co{int([1, 2])}).

\subsection{Ce qui vaut \co{False}}
\co{0}, \co{0.0}, \co{""}, \co{[]}, \co{\{\}}, \co{None}. Tout le reste est
vrai --- y compris \co{"0"}, \co{"False"} et \co{" "}.

% =====================================================================
\section{Opérateurs}

\begin{tabularx}{\linewidth}{@{}l>{\raggedright\arraybackslash}X@{}}
\toprule
\co{+ - *} & addition, soustraction, multiplication \\
\co{/}     & division \textbf{réelle} : toujours un \co{float} \\
\co{//}    & division entière, arrondie \textbf{vers le bas} \\
\co{\%}    & modulo (le reste) \\
\co{**}    & puissance, associative \textbf{à droite} \\
\bottomrule
\end{tabularx}

\vspace{0.4em}
Affectations combinées : \co{+=}, \co{-=}, \co{*=}, \co{/=}, \co{//=},
\co{\%=}, \co{**=}. \co{x += y} équivaut à \co{x = x + y}.

% =====================================================================
\section{Conditions}

\begin{lstlisting}
if a > 0:
    ...
elif a < 0:
    ...
else:
    ...
\end{lstlisting}

Comparaisons \textbf{chaînables} : \co{if 0 < a < 10:} équivaut à
\co{0 < a and a < 10}.

\subsection{\co{==} contre \co{is}}
\co{==} compare le \textbf{contenu}, \co{is} compare l'\textbf{identité} en
mémoire.

\begin{lstlisting}
a = [1, 2, 3]
b = [1, 2, 3]
c = a
a == b   # True
a is b   # False  (deux objets)
a is c   # True   (meme objet)
\end{lstlisting}

\subsection{Condition sur une ligne}
\co{res = "pair" if a \% 2 == 0 else "impair"}. Python n'a pas d'opérateur
ternaire \co{? :}.

\subsection{\co{in}}
Teste l'appartenance. Sur une chaîne, il cherche une \textbf{sous-chaîne} :
\co{"and" in "cat and dog"} vaut \co{True}.

% =====================================================================
\section{Boucles}

\begin{lstlisting}
range(5)        # 0 1 2 3 4
range(5, 10)    # 5 6 7 8 9
range(5, 10, 2) # 5 7 9
\end{lstlisting}

La borne de fin est \textbf{toujours exclue}.

\begin{itemize}
  \item \co{break} : sort immédiatement de la boucle
  \item \co{continue} : passe à l'itération suivante
  \item \co{else} : s'exécute \textbf{seulement} si la boucle s'est terminée
        \textbf{sans} \co{break}
\end{itemize}

\begin{lstlisting}
for i in range(10):
    if i == 5:
        break
else:
    print("fin")   # jamais affiche
\end{lstlisting}

% =====================================================================
\section{Listes}

\begin{tabularx}{\linewidth}{@{}l>{\raggedright\arraybackslash}X@{}}
\toprule
\co{len(l)} & la taille \\
\co{min/max/sum(l)} & agrégations \\
\co{sorted(l)} & \textbf{nouvelle} liste triée \\
\co{l.sort()} & trie \textbf{sur place}, renvoie \co{None} \\
\co{l.sort(reverse=True)} & ordre décroissant \\
\co{l.reverse()} & inverse l'ordre, ne trie pas \\
\co{l.append(x)} & ajoute \textbf{un} élément \\
\co{l.extend([x,y])} & ajoute \textbf{chaque} élément \\
\co{l.remove(v)} & supprime la 1re occurrence \\
\co{l.pop(i)} & supprime \textbf{et renvoie} l'index \co{i} \\
\co{l.pop()} & le dernier élément \\
\co{del l[i]} & supprime l'index \co{i} \\
\co{l.index(v)} & 1er index de \co{v} \\
\co{l.count(v)} & nombre d'occurrences \\
\co{l.clear()} & vide la liste \\
\bottomrule
\end{tabularx}

\vspace{0.4em}
Une liste est \textbf{hétérogène} : \co{[0, "test", 1.5, True]} est valide.

\subsection{Compréhension}
\co{[expression for e in iterable if condition]}

\begin{lstlisting}
[x ** 2 for x in range(10) if x % 2 != 0]
# [1, 9, 25, 49, 81]
\end{lstlisting}

\subsection{\co{enumerate} et \co{zip}}
\begin{lstlisting}
for i, fruit in enumerate(fruits):
    ...                  # index ET valeur

for nom, age in zip(noms, ages):
    ...                  # deux listes en parallele
\end{lstlisting}
\co{zip} s'arrête à la séquence la plus courte, sans erreur.

\subsection{Unpacking}
\co{prenom, nom = ['Homer', 'Simpson']} --- il faut autant de variables que
d'éléments.

% =====================================================================
\section{Slicing}

Syntaxe \co{[debut:fin\_exclue:pas]}. Produit une \textbf{copie}, jamais une
référence.

\begin{lstlisting}
n = [0, 1, 2, 3, 4, 5]
n[1:4]     # [1, 2, 3]
n[::2]     # [0, 2, 4]
n[::-1]    # [5, 4, 3, 2, 1, 0]  (copie inversee)
n[::-2]    # [5, 3, 1]
n[:]       # copie complete
n[1:4] = [10, 20]  # remplace la tranche
n[1:4] = []        # supprime la tranche
\end{lstlisting}

% =====================================================================
\section{Chaînes de caractères}

Les chaînes sont \textbf{immuables} : toute méthode renvoie une
\textbf{nouvelle} chaîne, il faut réaffecter.

\begin{lstlisting}
mot = mot.replace("l", "x")   # sinon aucun effet
\end{lstlisting}

\begin{tabularx}{\linewidth}{@{}l>{\raggedright\arraybackslash}X@{}}
\toprule
\co{.lower() .upper()} & la casse \\
\co{.split(sep)} & découpe, renvoie une \textbf{liste} \\
\co{.split()} & sépare sur espaces, \co{\textbackslash t}, \co{\textbackslash n} \\
\co{.split(" ", maxsplit=2)} & 2 découpages, donc 3 éléments \\
\co{sep.join(liste)} & assemble, renvoie une \textbf{chaîne} \\
\co{.strip()} & retire espaces et \co{\textbackslash n} aux \textbf{bords} \\
\co{.replace(a, b)} & remplace \\
\co{.find(s)} & l'index, ou \co{-1} si absent \\
\bottomrule
\end{tabularx}

% =====================================================================
\section{Dictionnaires}

\begin{lstlisting}
p = {"nom": "Alice", "age": 25}
p["age"] = 26        # modification
p["ville"] = "Paris" # ajout (meme syntaxe)

p.keys()    # les cles
p.values()  # les valeurs
p.items()   # les couples (cle, valeur)

for cle, valeur in p.items():
    print(cle, "->", valeur)
\end{lstlisting}

\subsection{\co{get} plutôt que \co{[]}}
\co{p["z"]} lève une \co{KeyError} ; \co{p.get("z")} renvoie \co{None} et
\co{p.get("z", 0)} renvoie \co{0}. D'où le comptage d'occurrences :

\begin{lstlisting}
compteur[mot] = compteur.get(mot, 0) + 1
\end{lstlisting}

\subsection{\co{setdefault} pour grouper}
\co{res.setdefault(cle, []).append(x)} crée la liste si besoin, puis ajoute.

\subsection{Compréhension}
\begin{lstlisting}
{x: x**2 for x in range(5)}
# {0: 0, 1: 1, 2: 4, 3: 9, 4: 16}

{c: v for c, v in d.items() if v % 2 == 0}
\end{lstlisting}

% =====================================================================
\section{Tuples et ensembles}

\textbf{Tuple} : séquence \textbf{immuable}. Modifier un élément lève une
\co{TypeError}. Il sert de clé de dictionnaire et de retour multiple.

\begin{lstlisting}
def min_max(l):
    return min(l), max(l)   # renvoie un TUPLE

x, y = min_max([1, 2, 3])   # unpacking
\end{lstlisting}

\textbf{Set} : collection \textbf{non ordonnée}, \textbf{sans doublon}.
\co{set([1, 1, 2, 2, 3])} donne \co{\{1, 2, 3\}}.

\subsection{Conversions entre structures}
\begin{lstlisting}
tuple(liste)          # figer -> utilisable comme cle
list(mon_tuple)       # pour modifier
dict([("a", 1)])      # liste de couples -> dict
dict(zip(cles, vals)) # deux listes -> dict
list(d.items())       # dict -> liste de tuples
\end{lstlisting}

Une \textbf{liste ne peut pas} être une clé de dictionnaire : elle est mutable,
donc non hachable, d'où une \co{TypeError}.

% =====================================================================
\section{Matrices}

Une matrice est une \textbf{liste de listes}. \co{len(m)} donne le nombre de
lignes, \co{len(m[0])} le nombre de colonnes. Les lignes peuvent avoir des
tailles différentes.

\begin{lstlisting}
[ligne[1] for ligne in m]     # la colonne 1
[list(c) for c in zip(*m)]    # la transposee
[[0 for _ in range(3)] for _ in range(3)]
\end{lstlisting}

\begin{piegebox}{Le piège de la duplication}
\co{[[0] * 3] * 3} duplique la \textbf{référence} : les trois lignes sont le
\textbf{même objet}. Modifier \co{m[0][0]} modifie les trois lignes. Toujours
passer par une compréhension.
\end{piegebox}

% =====================================================================
\section{Fonctions}

\begin{lstlisting}
def saluer(nom, message="Bonjour"):
    return f"{message}, {nom} !"

saluer("Bob")                   # defaut utilise
saluer("Bob", "Salut")          # positionnel
saluer(message="Hey", nom="Bo") # nomme
\end{lstlisting}

Sans \co{return}, une fonction renvoie \co{None}.

\subsection{Arguments variables}
\co{*args} donne un \textbf{tuple} d'arguments positionnels.
\co{**kwargs} donne un \textbf{dictionnaire} d'arguments nommés.

\subsection{Portée}
Une affectation dans une fonction crée une variable \textbf{locale} qui masque
la globale :

\begin{lstlisting}
x = 10
def f():
    x = 20      # variable LOCALE
f()
print(x)        # 10
\end{lstlisting}

\subsection{Annotations}
\co{def f(n: float) -> int:} n'impose rien. Aucune erreur si le type renvoyé
diffère.

% =====================================================================
\section{Lambda, map, filter, sorted}

\co{lambda arg1, arg2: expression} --- fonction anonyme d'une seule
expression, dont la valeur est renvoyée implicitement.

\begin{lstlisting}
carre = lambda x: x * x

list(map(lambda x: x**2, nombres))
list(filter(lambda x: x % 2 == 0, nombres))

sorted(eleves, key=lambda e: e["note"],
       reverse=True)
\end{lstlisting}

\co{map} et \co{filter} renvoient un \textbf{objet paresseux} : sans
\co{list(...)}, \co{print} affiche \co{<map object at ...>}. En pratique, la
compréhension de liste est préférée pour la lisibilité.

% =====================================================================
\section{Exceptions}

\begin{lstlisting}
try:
    ...
except ValueError:
    ...             # ce cas precis
except Exception as e:
    print(e)        # le fourre-tout, EN DERNIER
else:
    ...             # si AUCUNE exception
finally:
    ...             # TOUJOURS
\end{lstlisting}

\begin{tabularx}{\linewidth}{@{}l>{\raggedright\arraybackslash}X@{}}
\toprule
\textbf{Situation} & \textbf{Erreur} \\
\midrule
\co{int("abc")}            & \co{ValueError} \\
\co{int([1, 2])}           & \co{TypeError} \\
\co{l.remove(v)} absent    & \co{ValueError} \\
\co{l.index(v)} absent     & \co{ValueError} \\
\co{l.pop(i)} hors bornes  & \co{IndexError} \\
\co{d["z"]} absent         & \co{KeyError} \\
\co{10 / 0}                & \co{ZeroDivisionError} \\
\co{t[0] = x} sur un tuple & \co{TypeError} \\
liste comme clé de dict    & \co{TypeError} \\
\bottomrule
\end{tabularx}

\vspace{0.4em}
\co{raise ValueError("message")} lève volontairement une exception et
interrompt l'exécution.

% =====================================================================
\section{Affichage et fichiers}

\begin{lstlisting}
print(1, 2, 3)              # 1 2 3
print(1, 2, 3, sep="---")   # 1---2---3

print(f"Bonjour {nom}")
print(f"{1 + 2}")           # 3
print(f"{age:03d}")         # 025
print(f"{nom:>10}")         # cale a droite sur 10
\end{lstlisting}

\co{>} aligne à droite, \co{<} à gauche, \co{\textasciicircum} centre ; le
nombre est la largeur totale.

\begin{lstlisting}
with open("f.csv", mode="w",
          encoding="utf-8") as f:
    f.write(ligne + "\n")

with open("f.csv", mode="r",
          encoding="utf-8") as f:
    for ligne in f:            # economise la RAM
        nom, age = ligne.strip().split(",")
\end{lstlisting}

\co{with} \textbf{ferme automatiquement} le fichier, même en cas d'erreur.
Modes : \co{"r"} lecture, \co{"w"} écriture (\textbf{écrase tout}), \co{"a"}
ajout à la fin.

\subsection{Point d'entrée}
\co{if \_\_name\_\_ == "\_\_main\_\_":} n'exécute le bloc que si le fichier est
\textbf{lancé directement}, pas s'il est importé.

% =====================================================================
\section{Réflexes des exercices}

\begin{itemize}
  \item \textbf{Dédoublonner} : \co{set(l)} perd l'ordre,
        \co{list(dict.fromkeys(l))} le conserve.
  \item \textbf{Croissance stricte} :
        \co{all(a < b for a, b in zip(l, l[1:]))} compare chaque élément au
        suivant.
  \item \textbf{Court-circuit} : dans \co{len(m) > 5 and m[0] in VOYELLES}, le
        test de droite n'est évalué que si celui de gauche est vrai, ce qui
        évite une \co{IndexError} sur une chaîne vide.
  \item \textbf{Compter} : \co{c[mot] = c.get(mot, 0) + 1}, ou
        \co{Counter(mots).most\_common(2)}.
  \item \textbf{Inverser un dict} : \co{\{v: k for k, v in d.items()\}} --- si
        les valeurs ne sont pas uniques, la dernière écrase les précédentes.
  \item \textbf{Mémoïsation} : un dictionnaire sert de cache, chaque valeur
        n'est calculée qu'une seule fois (Fibonacci récursif).
  \item \textbf{Moyennes} : la moyenne d'une classe se calcule sur
        \textbf{toutes les notes} mises à plat, pas sur la moyenne des
        moyennes (sauf si chacun a le même nombre de notes).
  \item \textbf{Division par zéro} : tester \co{if not notes: return None}
        avant \co{sum(notes) / len(notes)}.
\end{itemize}

\end{multicols}

\vspace{0.3em}
\begin{astucebox}{Checklist des 60 dernières secondes}
\begin{minipage}[t]{0.485\linewidth}
\begin{itemize}
  \item \co{/} rend un \co{float}, \co{//} arrondit \textbf{vers le bas}.
  \item \co{.sort()} renvoie \co{None}, \co{sorted()} renvoie la liste.
  \item Chaînes et tuples sont \textbf{immuables} : il faut réaffecter.
  \item La borne de fin est \textbf{exclue} : \co{range} comme slicing.
\end{itemize}
\end{minipage}\hfill
\begin{minipage}[t]{0.485\linewidth}
\begin{itemize}
  \item \co{append} ajoute 1 élément, \co{extend} les ajoute tous.
  \item Le \co{else} d'une boucle veut dire \textbf{pas} de \co{break}.
  \item \co{*args} donne un tuple, \co{**kwargs} un dictionnaire.
  \item \co{finally} s'exécute \textbf{toujours}, \co{else} seulement si tout
        s'est bien passé.
\end{itemize}
\end{minipage}
\end{astucebox}

\vspace{0.4em}
\begin{center}
\footnotesize\itshape
Pour s'entraîner : \texttt{cd "Lesson 1" \&\& python quiz.py}
\end{center}

\end{document}
